\section{Positioning}
\label{sec:sota-positioning}

The four preceding sections leave four separate openings. Review-comment
generation is studied with the diff as input, and the benchmarks that enrich
context do not vary how it is selected. Review-comment evaluation has moved to
criterion-based and reference-free scoring but scores with a single instrument
whose reliability is not reported. The LLM-as-a-judge literature has developed the
countermeasures for exactly that problem but they have not been applied in this
setting. And repository graphs have been shown to help on tasks with checkable
outcomes, but review-comment generation remains unexamined.

Table~\ref{tab:sota-positioning} places this thesis against the four most directly
comparable systems: the two closest works on evaluating review comments, and the
two closest on graph-structured repository context.

\begin{table}[htbp]
    \centering
    \caption{Positioning against the closest prior work. \textbf{DC} =
             DeepCRCEval~\cite{hong2024deepcrceval} and \textbf{CR} =
             CRScore~\cite{naik2025crscore} are the closest work on
             \emph{evaluating} review comments; \textbf{GC} =
             GraphCoder~\cite{liu2024graphcoder} and \textbf{KGC} =
             KGCompass~\cite{yang2025kgcompass} are the closest work on
             \emph{graph-structured} repository context. ``n/a'' marks a dimension
             that does not apply because the system is scored by program execution
             rather than by judgement.}
    \label{tab:sota-positioning}
    \small
    \begin{tabularx}{\textwidth}{@{}Xccccc@{}}
        \toprule
          & DC & CR & GC & KGC & \textbf{This thesis} \\
        \midrule
        Task is review comments
          & $\bullet$ & $\bullet$ & --- & --- & $\bullet$ \\
        Adds out-of-diff context
          & --- & --- & $\bullet$ & $\bullet$ & $\bullet$ \\
        Structural (graph) context
          & --- & --- & $\bullet$ & $\bullet$ & $\bullet$ \\
        Retrieval and structure as matched arms
          & --- & --- & $\bullet$ & --- & $\bullet$ \\
        Independent LLM scorers
          & 1 & 1 & n/a & n/a & \textbf{5} \\
        Inter-scorer agreement reported
          & --- & --- & n/a & n/a & $\bullet$ \\
        Detection adjudication checked against oracle
          & --- & --- & n/a & n/a & $\bullet$ \\
        Human corroboration of targeted criterion
          & $\bullet$ & $\bullet$ & --- & --- & $\bullet$ \\
        \bottomrule
    \end{tabularx}
\end{table}

One row deserves comment, because it is the row on which this thesis is least
novel. GraphCoder does compare structural retrieval against similarity retrieval
as matched arms, evaluating its code-context-graph retrieval against vanilla,
shifted and iterative sequence-based retrieval with the same models on the same
completion tasks~\cite{liu2024graphcoder}. The comparison this thesis makes is
therefore not new as a study design; what is new is the task it is applied to and
the instrument used to read it. That distinction narrows the claim below to
what the evidence supports.

\subsection*{Research Gap}

No existing study compares structural and similarity-based out-of-diff context
for generating code-review comments within one matched generation and
evaluation pipeline, while also measuring sensitivity to the evaluator.
GraphCoder runs a matched comparison on a task with a checkable outcome, while
DeepCRCEval and CRScore evaluate review comments but do not vary context or
report scorer sensitivity.

This thesis occupies that position along four dimensions that the table makes
visible.

It compares selection strategies rather than proposing one. Structural context and
similarity retrieval are supplied to the same model, on the same pull requests,
inside a matched prompt template, and are also supplied together. The arms vary
both the selected material and the instruction sentence that names it; this
confound is described in Chapter~\ref{chap:approach} and treated explicitly in
Chapter~\ref{chap:results}.

It treats the evaluation instrument as an object of study. The rubric is scored by
a panel drawn from more than one provider rather than by a single judge.
Chapter~\ref{chap:results} reports how much the conclusion moves across
individual judges and generators, a sensitivity analysis that the
single-judge protocols cannot perform.

Its structural context is derived from a code property graph with resolved
call and data-flow edges; the practitioner study additionally evaluates a
lightweight import resolver to assess whether the benefit holds under reduced
analysis depth.

It corroborates a coverage rubric with two instruments that can see what coverage
cannot. A controlled defect-injection experiment scores detection against a
deterministic oracle rather than a judge, and a human study asks practitioners
which review they would rather receive. Where the three instruments
agree, the finding is stronger than any of them alone; where they disagree,
Chapter~\ref{chap:results} says so.
